\documentclass[11pt]{article}
\usepackage[margin=1in]{geometry}
\usepackage{amsmath}
\usepackage{amssymb}
\usepackage{array}
\usepackage{booktabs}
\usepackage{longtable}
\usepackage{hyperref}

\title{QMM Complete Workflow Manual}
\author{Diego Juarez and Codex}
\date{\today}

\begin{document}
\maketitle
\tableofcontents
\newpage

\section{Purpose}

This document explains the workflow implemented in \texttt{QMM}.  The goal is
to keep the physics, numerical methods, code organization, and JSON interface
in one clean place.

The package is designed to support two use cases:

\begin{itemize}
\item production runs for the models already implemented in
\texttt{src/qmm/models.py}
\item extension to a new real-gas model with a new mean field and/or a new
excluded-volume prescription
\end{itemize}

\section{QMM Organization}

The main directories are:

\begin{itemize}
\item \texttt{src/qmm/}: source code
\item \texttt{examples/}: JSON inputs, including a blank generalized template
\item \texttt{notebooks/QMM\_workflow.ipynb}: single overview notebook
\item \texttt{docs/qmm\_complete\_workflow\_manual.tex}: this manual
\item \texttt{results/}: local reference outputs used by the notebook
\end{itemize}

The main solver modules are:

\begin{itemize}
\item \texttt{models.py}: interaction models and repulsive mappings
\item \texttt{fermi\_gas.py}: relativistic ideal Fermi gas
\item \texttt{ground\_state.py}: symmetric zero-temperature fit
\item \texttt{quantum.py}: finite-temperature liquid critical point
\item \texttt{quarkyonic.py}: symmetric and asymmetric quarkyonic matter
\item \texttt{sound\_speed.py}: sound-speed reconstruction
\item \texttt{asymmetry.py}: asymmetric hadronic couplings
\item \texttt{neutron\_star.py}: beta-equilibrium EOS preparation and TOV solver
\item \texttt{runner.py}: top-level workflow orchestration
\item \texttt{config.py}: JSON loading and typed configuration
\end{itemize}

\section{What The JSON File Controls}

The workflow is driven by a single JSON input file.  The blank template is:

\begin{center}
\texttt{examples/qmm\_blank\_template.json}
\end{center}

The richer single-file design template is:

\begin{center}
\texttt{examples/qmm\_all\_in\_one\_template.json}
\end{center}

The JSON file can decide:

\begin{itemize}
\item which registered model is used
\item whether an extra model parameter is fixed or solved from a target
incompressibility
\item whether to compute the symmetric ground-state point
\item whether to compute the liquid critical point
\item whether to compute a dense pure-hadronic EOS table
\item whether to compute symmetric quarkyonic matter
\item whether to compute asymmetric couplings
\item whether to compute asymmetric fixed-$y$ quarkyonic matter
\item whether to include leptons through beta equilibrium
\item whether to compute neutron stars
\item whether CSV, JSON, and plot outputs are written
\end{itemize}

\subsection{Documentation Blocks Inside The Same JSON}

The all-in-one template also contains descriptive blocks such as:

\begin{itemize}
\item \texttt{documentation.mean\_field\_latex}
\item \texttt{documentation.excluded\_volume\_latex}
\item \texttt{documentation.ideal\_density\_map\_latex}
\item \texttt{documentation.pressure\_prefactor\_latex}
\item \texttt{requested\_calculations.*}
\end{itemize}

These fields are intentionally not executable physics inputs.  They are there
so one JSON file can also serve as a human-readable model record.  The
\texttt{QMM} loader ignores these extra descriptive keys.

\subsection{What The JSON Cannot Do By Itself}

The JSON file does not encode new analytic physics functions.  In particular,
it cannot by itself define:

\begin{itemize}
\item a new mean field $U(n)$
\item a new derivative $dU/dn$
\item a new density mapping $n \leftrightarrow n_{\mathrm{id}}$
\item a new excluded-volume fraction
\item a new pressure prefactor
\end{itemize}

Those must be implemented once in \texttt{src/qmm/models.py}.  After the model
is registered there, the same JSON interface can drive it.

\section{Units and Default Physical Inputs}

The code uses MeV--fm units for the nuclear-matter solvers and converts to cgs
only when solving the TOV equations.

The default physical constants are:

\begin{align}
N_c &= 3  \quad \text{color}, \\
g &= 4 \quad \text{for symmetric nuclear matter}, \\
g &= 2 \quad \text{for each nucleon species in asymmetric matter}, \\
\hbar c &= 197.3269804~\text{MeV fm}, \\
m_p &= 938.27208816~\text{MeV}, \\
m_N &= 938.0~\text{MeV}, \\
m_u^{\overline{\mathrm{MS}}}(2~\text{GeV}) &= 2.16~\text{MeV}, \\
m_d^{\overline{\mathrm{MS}}}(2~\text{GeV}) &= 4.67~\text{MeV}, \\
n_0 &= 0.16~\text{fm}^{-3}, \\
E/A - m_N &= -16~\text{MeV at saturation}, \\
J &= 32.5~\text{MeV}, \\
L &= 58.9~\text{MeV}.
\end{align}


\section{Interaction Models}

Each model is represented by one \texttt{InteractionModel} entry in
\texttt{models.py}.  The definition requires the following ingredients:

\begin{itemize}
\item attractive function $U(n,b,\mathrm{parameter})$
\item derivative $dU/dn$
\item map $n \mapsto n_{\mathrm{id}}$
\item inverse map $n_{\mathrm{id}} \mapsto n$
\item total volume fraction
\item species-level volume fraction
\item pressure prefactor
\item capability flags for each workflow branch
\end{itemize}

\subsection{Attractive Mean Fields}

The attractive sectors implemented in \texttt{QMM} are:

\begin{align}
U_{\mathrm{vdW}}(n) &= -n, \\
U_{\mathrm{RKS}}(n) &= -\frac{1}{b}\ln(1 + bn), \\
U_{\mathrm{PR}}(n) &=
-\frac{1}{2\sqrt{2}\,b}
\ln \left[
\frac{1 + (1+\sqrt{2})bn}{1 + (1-\sqrt{2})bn}
\right], \\
U_{\mathrm{Clausius}}(n) &= -\frac{n}{1 + cn}, \\
U_{\mathrm{Dieterici}}(n) &= -\frac{n^{\alpha - 1}}{\alpha - 1}.
\end{align}

For the \texttt{cs} and \texttt{tvm} models, the attractive sector remains the
van der Waals form $U(n)=-n$, while the repulsive sector is modified through
the excluded-volume mapping.  The hybrid models \texttt{clausius\_cs} and
\texttt{clausius\_tvm} keep the Clausius attractive sector and only replace
the repulsive mapping.

\subsection{Excluded-Volume Prescriptions}

The repulsive sector in \texttt{QMM} is encoded through:

\begin{itemize}
\item a volume fraction $f(n,b)$
\item an ideal-density map $n \mapsto n_{\mathrm{id}}$
\item a pressure prefactor $\Pi(n,b)$
\end{itemize}

\subsubsection{vdW / Standard Excluded Volume}

The standard excluded-volume prescription used by
\texttt{vdw}, \texttt{rks}, \texttt{pr}, \texttt{clausius}, and
\texttt{dieterici} is:

\begin{align}
f_{\mathrm{EV}}(n) &= 1 - bn, \\
n_{\mathrm{id}} &= \frac{n}{1 - bn}, \\
\Pi_{\mathrm{EV}}(n,b) &= 1.
\end{align}

\subsubsection{Carnahan--Starling Excluded Volume}

For the \texttt{cs} model, and also for the hybrid \texttt{clausius\_cs},
define $x = bn$.  The excluded-volume fraction is:

\begin{align}
f_{\mathrm{CS}}(x) &=
\exp\left[
-\frac{3x}{4-x} - \frac{4x}{(4-x)^2}
\right], \\
n_{\mathrm{id}} &= \frac{n}{f_{\mathrm{CS}}(x)}.
\end{align}

The pressure prefactor is:

\begin{equation}
\Pi_{\mathrm{CS}}(n,b) =
f_{\mathrm{CS}}(n,b) - n \frac{d f_{\mathrm{CS}}(n,b)}{dn}.
\end{equation}

\subsubsection{TVM Excluded Volume}

For the \texttt{tvm} model, and also for the hybrid \texttt{clausius\_tvm}
model, define again $x = bn$.  The TVM excluded-volume fraction is:

\begin{align}
f_{\mathrm{TVM}}(x) &= \exp\left(-x - \frac{x^2}{2}\right), \\
n_{\mathrm{id}} &= \frac{n}{f_{\mathrm{TVM}}(x)}.
\end{align}

The pressure prefactor is:

\begin{equation}
\Pi_{\mathrm{TVM}}(n,b) =
f_{\mathrm{TVM}}(n,b) - n \frac{d f_{\mathrm{TVM}}(n,b)}{dn}.
\end{equation}

\subsubsection{Species-Level Excluded Volume in Asymmetric Matter}

For asymmetric matter, the code uses the species-dependent arguments

\begin{align}
x_p &= b_n n_p + b_{pn} n_n, \\
x_n &= b_{pn} n_p + b_n n_n,
\end{align}

and then evaluates the species volume fraction with the corresponding model
rule: $1-x$ for standard excluded volume, $f_{\mathrm{CS}}(x)$ for
Carnahan--Starling, or $f_{\mathrm{TVM}}(x)$ for TVM.

\subsection{Examples Already Implemented}

The package already contains the following models:

\begin{itemize}
\item \texttt{vdw}
\item \texttt{rks}
\item \texttt{pr}
\item \texttt{clausius}
\item \texttt{clausius\_cs}
\item \texttt{clausius\_tvm}
\item \texttt{dieterici}
\item \texttt{cs}
\item \texttt{tvm}
\end{itemize}

The hybrid model \texttt{clausius\_cs} combines the Clausius attractive sector
with the Carnahan--Starling excluded-volume mapping, while
\texttt{clausius\_tvm} combines the same Clausius attraction with the TVM
repulsive mapping.  Because these are not the original pure Clausius model,
the practical search window for the parameter $c$ may differ from the usual
pure-Clausius reference range.  In the included scan examples, the search
interval is widened to reach low-$K_0$ targets such as
$K_0=240~\mathrm{MeV}$.

\section{Workflow 1: Symmetric Zero-Temperature Ground State}

At saturation density the code solves for the couplings $a$ and $b$ using:

\begin{align}
P(n_0) &= 0, \\
\frac{\epsilon(n_0)}{n_0} - m_N &= -16~\mathrm{MeV}.
\end{align}

After solving $a$ and $b$, the incompressibility is evaluated as
\begin{equation}
K_0 = 9 \left. \frac{dP}{dn_B} \right|_{n_0}.
\end{equation}

The symmetric energy density and pressure use the model-dependent ideal-density
map and mean field:
\begin{align}
\epsilon(n) &= f(n,b)\,\epsilon_{\mathrm{id}}(n_{\mathrm{id}}) + n\,a\,U(n), \\
P(n) &= \Pi(n,b)\,P_{\mathrm{id}}(n_{\mathrm{id}}) + n^2 a \frac{dU}{dn}.
\end{align}

\section{Workflow 2: Finite-Temperature Liquid Critical Point}

The finite-temperature critical point is obtained from the interacting pressure
surface by solving:
\begin{align}
\left(\frac{\partial P}{\partial n_B}\right)_T &= 0, \\
\left(\frac{\partial^2 P}{\partial n_B^2}\right)_T &= 0.
\end{align}

The code first inverts the finite-temperature ideal-gas density to obtain the
effective chemical potential $\mu^\ast$, then reconstructs the interacting
pressure.  The outputs are:

\begin{itemize}
\item $T_c$
\item $n_c$
\item $P_c$
\end{itemize}

\section{Workflow 3: Symmetric Quarkyonic Matter}

At fixed baryon density $n_B$, the code minimizes the zero-temperature energy
density with respect to the quark fraction $f_Q$:
\begin{equation}
\epsilon = \epsilon(n_B, f_Q).
\end{equation}

The quarkyonic shell scale is controlled by the JSON entry
\texttt{quarkyonic.lambda\_momentum\_mev}.  This is the main place where the
CS/TVM versus other-model choices can differ in practice.

\subsection{Sound Speed}

Once the minimum-energy curve $\epsilon(n_B)$ has been generated, the code
reconstructs:
\begin{align}
\mu_B &= \frac{d\epsilon}{dn_B}, \\
P &= n_B \mu_B - \epsilon, \\
v_s^2 &= \frac{dP/dn_B}{\mu_B}.
\end{align}

In the current design, if the quarkyonic profile is requested, sound speed is
reconstructed automatically as part of the workflow.

\section{Workflow 3A: Pure Hadronic EOS Table}

The workflow can also write a dense zero-temperature hadronic EOS table before
the quarkyonic construction is applied.  This is controlled by:

\begin{center}
\texttt{workflows.hadronic\_eos\_table = true}
\end{center}

The density range and derivative-reconstruction settings are controlled by the
\texttt{hadronic\_eos} block.  The output table contains:

\begin{itemize}
\item $n$
\item $n/n_0$
\item raw and smoothed $\epsilon(n)$
\item $E/A - m_N$
\item $\mu_B$
\item $P$
\item $dP/dn_B$
\item $d^2 \epsilon / dn_B^2$
\item $v_s^2$
\end{itemize}

This provides a direct hadronic EOS and sound-speed diagnostic without turning
on the quarkyonic branch.

\section{Workflow 4: Asymmetric Hadronic Matter}

The asymmetric hadronic sector is fitted from the five constraints used in this
project:

\begin{itemize}
\item binding energy
\item saturation density
\item symmetry energy $J$
\item symmetry-energy slope $L$
\item incompressibility $K_0$
\end{itemize}

The asymmetric couplings are:
\begin{equation}
(a_n, a_{pn}, b_n, b_{pn}).
\end{equation}

The code currently supports two branches:

\begin{itemize}
\item \texttt{equal\_b}: $b_{pn} = b_n$
\item \texttt{target\_l}: $b_{pn} \neq b_n$, chosen to reproduce $L$
\end{itemize}

\section{Workflow 5: Asymmetric Quarkyonic Matter}

Once the asymmetric couplings are known, the code can compute:

\begin{itemize}
\item fixed-$y$ profiles for user-selected proton fractions
\item beta-equilibrium matter with leptons
\end{itemize}

The proton-fraction grid is selected in:
\begin{center}
\texttt{asymmetric.proton\_fraction\_values}
\end{center}

Leptons are controlled by:
\begin{center}
\texttt{asymmetric.beta\_equilibrium}
\end{center}

If \texttt{beta\_equilibrium=false}, the code computes only fixed-$y$ matter.
If \texttt{beta\_equilibrium=true}, the code solves for charge neutrality with
electrons and muons and constructs the beta-equilibrium EOS.

\section{Workflow 6: Neutron Stars}

The neutron-star branch uses the beta-equilibrium EOS and integrates the
Tolman--Oppenheimer--Volkoff equations with an RK4 solver.

The top-level switch is:
\begin{center}
\texttt{workflows.neutron\_star}
\end{center}

The main numerical controls are:

\begin{itemize}
\item number of central-pressure points
\item pressure range
\item radial step size
\item maximum radius
\end{itemize}

\section{How To Read The Workflow Flags}

\subsection{Minimal Symmetric Calibration}

If you only want to solve the symmetric ground-state point and obtain
$(a,b,K_0)$, set:
\begin{itemize}
\item \texttt{ground\_state=true}
\item everything else in \texttt{workflows} to \texttt{false}
\end{itemize}

\subsection{Add The Liquid Critical Point}

Set:
\begin{itemize}
\item \texttt{quantum\_critical=true}
\end{itemize}

\subsection{Add Symmetric Quarkyonic Matter}

Set:
\begin{itemize}
\item \texttt{symmetric\_quarkyonic=true}
\end{itemize}

\subsection{Add A Pure Hadronic EOS Table}

Set:
\begin{itemize}
\item \texttt{hadronic\_eos\_table=true}
\end{itemize}

\subsection{Add Asymmetric Matter Without Leptons}

Set:
\begin{itemize}
\item \texttt{asymmetric\_fit=true}
\item \texttt{asymmetric\_quarkyonic=true}
\item \texttt{asymmetric.beta\_equilibrium=false}
\end{itemize}

\subsection{Add Leptons}

Set:
\begin{itemize}
\item \texttt{asymmetric\_fit=true}
\item \texttt{asymmetric\_quarkyonic=true}
\item \texttt{asymmetric.beta\_equilibrium=true}
\end{itemize}

\subsection{Add Neutron Stars}

Set:
\begin{itemize}
\item \texttt{neutron\_star=true}
\end{itemize}

\section{How To Add A New Model}

For a new real-gas model, implement the model in \texttt{src/qmm/models.py}:

\begin{enumerate}
\item define the attractive sector $U(n)$ and $dU/dn$
\item define the excluded-volume mapping
\item define the inverse density mapping
\item define the volume fraction
\item define the pressure prefactor
\item register a new \texttt{InteractionModel}
\item set the capability flags for supported workflow branches
\end{enumerate}

After that, put the registered name into:
\begin{center}
\texttt{model.name}
\end{center}

inside a JSON file based on either \texttt{qmm\_blank\_template.json} or
\texttt{qmm\_all\_in\_one\_template.json}.

\section{Output Controls}

The \texttt{output} block controls what is written:

\begin{itemize}
\item \texttt{write\_csv}: write EOS tables and summaries
\item \texttt{write\_json}: write the run summary
\item \texttt{write\_plots}: write figures
\item \texttt{plot\_formats}: choose \texttt{png}, \texttt{pdf}, or both
\end{itemize}

\section{Practical Limitation}

In the current implementation, sound-speed reconstruction is tied to the
quarkyonic profile generation.  Therefore:

\begin{itemize}
\item if you request a quarkyonic branch, $v_s^2$ is reconstructed
\item if you do not want that branch, switch it off in \texttt{workflows}
\item if you want the branch but do not want saved files, switch off
\texttt{output.write\_csv} and/or \texttt{output.write\_plots}
\end{itemize}

\section{Main Files To Use}

For normal use, the main files are:

\begin{itemize}
\item \texttt{examples/qmm\_blank\_template.json}
\item \texttt{examples/qmm\_all\_in\_one\_template.json}
\item \texttt{notebooks/QMM\_workflow.ipynb}
\item \texttt{src/qmm/models.py}
\item \texttt{src/qmm/config.py}
\item \texttt{src/qmm/runner.py}
\end{itemize}

\end{document}
